%% parts/p07/structural-extensions.tex --- Phase-4 P7 new-content slice.
%% Coverage: N-a-056, N-a-057, N-a-058, D-P10-031; exclusions N-x-011/012.
%% No v1.1 proof status is changed: this chapter is register-driven content.

\chapter{Epsilon Calculus and Generalized Sequent Frameworks}
\label{sec:p07-structural-extensions}
\chaptermark{Epsilon and Generalized Frameworks}

\section{The equality-free epsilon calculus}
\label{sec:p07-epsilon-calculus}

\begin{hbdefinition}{def:p07-epsilon-calculus}{Equality-free epsilon calculus}
Fix an equality-free first-order language and keep its quantifiers primitive.
The calculus \(\mathsf{PC}_{\varepsilon}\) consists of all axioms and rules on
\hyperref[sheet:hilbert-fol]{\texttt{sheet:hilbert-fol}}, over the language
enlarged by the term former
\(\varepsilon x\,\varphi(x)\), together with the critical-axiom scheme on
\hyperref[sheet:epsilon]{\texttt{sheet:epsilon}}.  Formation and substitution
are capture-avoiding, and bound-variable variants name the same scheme.

The name is needed to separate this proof-theoretic witness term from the
set-theoretic Axiom of Choice and from the ordinal \(\varepsilon_0\).  Its
standing is a term-forming extension of one fixed equality-free calculus, not
a new semantic quantifier.  It is nonredundant because the critical axiom can
replace a witnessed formula by one selected term while ordinary predicate
syntax has no such term former.  The distinction is used by all three
conservativity statements below; equality-bearing epsilon calculi fall
outside their scope.
\index{epsilon calculus}\index{epsilon term}\index{critical axiom}
\end{hbdefinition}

Let \(\mathsf{EC}\) denote the equality-free, quantifier-free fragment of the
same classical calculus.  This metalinguistic abbreviation is needed only to
state the first epsilon theorem without repeatedly naming the fragment; it is
not a second epsilon system, and the theorem below is its later use.

\begin{hbtheorem}{thm:p07-first-epsilon}{cited=MoserZach2006}{First epsilon theorem}
Let \(T\) be a set of epsilon-free, quantifier-free formulas and let
\(\chi\) be epsilon-free and quantifier-free.  If \(\chi\) is derivable from \(T\)
in \(\mathsf{PC}_{\varepsilon}\), then \(\chi\) is derivable from \(T\) already
in \(\mathsf{EC}\).  The theorem here is equality-free
\cite{HilbertBernays1939,Buss1998,MoserZach2006}.
\end{hbtheorem}

\begin{hbtheorem}{thm:p07-extended-first-epsilon}{cited=MoserZach2006}{Extended first epsilon theorem and Herbrand disjunction}
Suppose the equality-free sentence
\[
  \exists x_1\cdots\exists x_k\,\varphi(x_1,\ldots,x_k)
\]
has quantifier-free matrix and is derivable in
\(\mathsf{PC}_{\varepsilon}\).  Then there are finitely many tuples of
epsilon-free terms \(\vec t_1,\ldots,\vec t_m\) such that the quantifier-free
classical calculus derives the Herbrand disjunction
\[
  \varphi(\vec t_1)\lor\cdots\lor\varphi(\vec t_m).
\]
This is the general-form bridge to the validity-form Herbrand theorem in
Part~2; it does not replace that part's universe, structure, or clause-set
entries \cite{Herbrand1971,Buss1998,MoserZach2006}.
\end{hbtheorem}

\begin{hbtheorem}{thm:p07-second-epsilon}{cited=MoserZach2006}{Second epsilon theorem}
Let \(T\) be a set of epsilon-free formulas of the equality-free first-order
language, and let \(\chi\) be an epsilon-free formula of that language.  If
\(\chi\) is derivable from \(T\) in
\(\mathsf{PC}_{\varepsilon}\), then \(\chi\) is derivable from \(T\) in
ordinary equality-free classical predicate calculus.  No equality version is
inferred from this statement.
\end{hbtheorem}

The three theorem boundaries, including the equality-free restriction, are
documented in \textcite{MoserZach2006}.

\begin{hbwarning}{def:p07-epsilon-boundary}{What the epsilon theorems do not say}
The first theorem is quantifier-free conservativity, the extended first
theorem extracts a finite Herbrand disjunction, and the second theorem removes
epsilon terms from an ordinary first-order consequence.  They are not one
interchangeable claim, and none supplies the more involved equality case.
\end{hbwarning}

\section{Four generalized sequent architectures}
\label{sec:p07-generalized-sequent-frameworks}

\begin{hbdefinition}{def:p07-generalized-sequent-frameworks}{Generalized sequent frameworks}
The following four names are required because they alter different parts of
a derivation object; calling all of them ``nonstandard sequents'' would erase
the distinction needed to choose a rule sheet.
\begin{enumerate}
  \item A \emph{hypersequent} is a finite sequence or multiset of sequents,
        with a separately declared external separator and separately declared
        external structural rules.
  \item A \emph{labelled sequent} contains labelled formulas such as
        \(w{:}\varphi\) and relational atoms such as \(wRv\), so a selected
        Kripke frame relation occurs inside the derivation language.
  \item A \emph{display calculus} has a structural language and reversible
        display postulates with the display property: any chosen substructure
        of a sequent can be made the entire antecedent or entire consequent of
        a structurally equivalent sequent.
  \item A \emph{deep-inference calculus} permits a rule instance inside an
        arbitrary formula or structure context, rather than only at its root.
\end{enumerate}
Their standing is architectural: each describes a class of derivation
objects, not a single logic or an adequacy theorem.  They are nonredundant
  respectively because they add external components, semantic labels, a
  residuated structural language, or non-root rule application.  The comparison
  table and the selected system-specific examples below use exactly these four
  distinctions
  \cite{Avron1996,Negri2005,Belnap1982,Wansing1994,Guglielmi2007,Stouppa2007}.
\index{hypersequent}\index{labelled sequent}\index{display calculus}
\index{deep inference}
\end{hbdefinition}

\begingroup\small
\begin{longtable}{@{}p{0.17\textwidth}p{0.28\textwidth}p{0.23\textwidth}p{0.22\textwidth}@{}}
\caption{Generalized sequent architectures and the selected system-specific
  modal evidence used by the worked examples.}
\label{tab:p07-generalized-frameworks}\\
\toprule
Architecture & Derivation object & Distinguishing mechanism & Named modal evidence boundary \\
\midrule
\endfirsthead
\toprule
Architecture & Derivation object & Distinguishing mechanism & Named modal evidence boundary \\
\midrule
\endhead
Hypersequent & Finite family of ordinary sequents & External separator and external structural rules & Avron's fixed GS5 presentation; no family-wide rule profile is inferred. \\
\addlinespace
Labelled & Labelled formulas plus relational atoms & Frame relation and geometric frame rules occur syntactically & Negri's labelled K calculus; T, S4, and S5 require separate frame rules. \\
\addlinespace
Display & Sequents in a structural language & Display postulates isolate a chosen substructure & A named Wansing modal display calculus, not Belnap's generic property alone. \\
\addlinespace
Deep inference & Formulas or structures with contexts at every depth & Rules may fire below the root & Stouppa's named S5 system; Guglielmi's BV supplies only the general architecture. \\
\bottomrule
\end{longtable}
\endgroup

\begin{hbwarning}{def:p07-worked-derivation-boundary}{A worked derivation is not adequacy}
Even after a system-specific sheet exists, one complete derivation proves only
that its endformula is derivable.  Soundness, completeness, cut admissibility,
and correspondence with a frame class require separate status-bearing
results.
\end{hbwarning}

\subsection*{Four matched modal derivations}

\begin{hbexample}{def:p07-gs5-derivation}{Avron GS5 derives axiom 5}
On \hyperref[sheet:gs5-avron]{\texttt{sheet:gs5-avron}}, start with
\(\Box p\Rightarrow\Box p\). Lifted negation-left gives
\(\Box p,\neg\Box p\Rightarrow\). Modalized Splitting gives
\[
 \Box p\Rightarrow\ \mid\ \neg\Box p\Rightarrow.
\]
Lifted negation-right in the first component yields
\(\Rightarrow\neg\Box p\mid\neg\Box p\Rightarrow\), and Box-R yields
\(\Rightarrow\Box\neg\Box p\mid\neg\Box p\Rightarrow\).
Internal weakening makes both components
\(\neg\Box p\Rightarrow\Box\neg\Box p\); external contraction merges them,
and implication-right concludes
\(\Rightarrow\neg\Box p\to\Box\neg\Box p\). Every structural step is on the
sheet, and the lifted propositional rules are those of classical LK
\cite{Avron1996}.
\end{hbexample}

\begin{hbexample}{def:p07-labelled-k-derivation}{Negri labelled \(K\) derives the K axiom}
Work backward on
\(\Rightarrow x:\Box(A\to B)\to(\Box A\to\Box B)\) using
\hyperref[sheet:labelled-k-negri]{\texttt{sheet:labelled-k-negri}}. Two
Implication-R steps reduce the goal to
\(x:\Box(A\to B),x:\Box A\Rightarrow x:\Box B\). Box-R introduces a fresh
\(y\) and relational atom \(xRy\). Two Box-L steps expose
\(y:A\to B\) and \(y:A\). Implication-L closes the remaining goal from the
identity sequents \(y:A\Rightarrow y:A\) and
\(y:B\Rightarrow y:B\). The freshness of \(y\) is the only eigenlabel
condition; no reflexivity or transitivity rule is used, so the derivation is in
basic \(K\) \cite{Negri2005}.
\end{hbexample}

\begin{hbexample}{def:p07-display-k-derivation}{A complete modal-display cycle in Wansing's calculus}
On \hyperref[sheet:display-k-wansing]{\texttt{sheet:display-k-wansing}}, the
identity \(A\Rightarrow A\) first gives
\(\Box A\Rightarrow\bullet A\) by Box-L.  Apply Modal Display in its declared
reverse direction to obtain
\(\bullet\Box A\Rightarrow A\), then apply Box-R:
\[
 A\Rightarrow A
 \;\overset{\Box L}{\Longrightarrow}\;
 \Box A\Rightarrow\bullet A
 \;\overset{\text{Modal Display}^{-1}}{\Longrightarrow}\;
 \bullet\Box A\Rightarrow A
 \;\overset{\Box R}{\Longrightarrow}\;
 \Box A\Rightarrow\Box A.
\]
Every intermediate sequent is therefore reconstructible from the three named
rows of the printed sheet.  The example exhibits the structural modal operator
on both sides of a display sequent; it is not offered as a derivation of the K
axiom or as an adequacy proof \cite{Wansing1994}.
\end{hbexample}

\begin{hbexample}{def:p07-deep-s5-derivation}{Stouppa deep inference derives axiom B}
On \hyperref[sheet:deep-s5-stouppa]{\texttt{sheet:deep-s5-stouppa}}, apply
i-down to \(\mathsf{tt}\) with structure \(\Diamond\Box a\), obtaining
\([\Diamond\Box a,\Box\Diamond\bar a]\). Apply 5-down inside the left
disjunct to obtain \([\Box a,\Box\Diamond\bar a]\), then t-down there to obtain
\([a,\Box\Diamond\bar a]\). With \(a=\neg p\), this is
\(p\to\Box\Diamond p\), axiom B. Both modal rules act strictly below the outer
disjunction, which is the promised deep-inference use \cite{Stouppa2007}.
\end{hbexample}

\begin{hbtheorem}{thm:p07-gs5-adequacy}{cited=Avron1996}{Adequacy and cut for the selected GS5 hypersequent calculus}
The selected GS5 presentation is sound and complete for \(\mathsf{S5}\), and
its cited cut-elimination theorem applies to that presentation.
\end{hbtheorem}

\begin{hbtheorem}{thm:p07-labelled-k-adequacy}{cited=Negri2005}{Adequacy and cut for the selected labelled K calculus}
The selected labelled calculus is sound and complete for basic normal modal
\(\mathsf K\), and Cut is admissible under its cited labelled-calculus
metatheorem.
\end{hbtheorem}

\begin{hbtheorem}{thm:p07-display-k-adequacy}{cited=Wansing1994}{Adequacy and cut for Wansing's full display K calculus}
Wansing's full display calculus is sound and complete for basic normal modal
\(\mathsf K\), and Cut is admissible under the source's display-calculus
metatheorem.  The printed sheet is only the worked-derivation fragment and is
not asserted to inherit those metatheorems as a complete standalone calculus.
\end{hbtheorem}

\begin{hbtheorem}{thm:p07-deep-s5-adequacy}{cited=Stouppa2007}{Adequacy and cut for the selected deep-inference S5 calculus}
The selected \(\mathsf{KSg\mbox{-}S5}\) presentation is sound and complete for
\(\mathsf{S5}\), and it has the cut-admissibility result stated in the cited
presentation.
\end{hbtheorem}

The four worked derivations prove only their displayed endformulas.  The four
theorems above keep each adequacy and cut claim attached to the source that
actually establishes it.

The cross-domain survey is now explicit in
Table~\ref{tab:p07-generalized-frameworks}. Hypersequents add an external
component layer; labelled calculi internalize relational semantics; display
calculi add a residuated structural language; deep inference changes the
allowed rule depth. Part VIII supplies the modal, intuitionistic, relevance,
and linear systems to which particular instances apply. No row in the table
asserts family-wide soundness, completeness, or cut elimination: those remain
properties of a named calculus with a named semantics.

\begin{proofboundary}{def:p07-proof-mining-pointer}{Proof-mining pointer}
Proof mining applies interpretations such as monotone functional
interpretation, majorization, and bounded interpretations to a formalized proof
in order to extract an explicit numerical bound or rate whose dependencies on
parameters can be audited. The term is needed to distinguish quantitative
extraction from ordinary program extraction or proof search. It stands as a
metamathematical analysis of a specified proof, differs because the output is
a bound with tracked dependence rather than a proof term for the original
statement, and is used here only as the route to the specialist development
\cite{Kohlenbach2008}. No uniform bound is claimed without the theorem form,
base theory, interpretation, and majorizability hypotheses required there.
\index{proof mining}
\end{proofboundary}

Geometry of interaction and proof nets remain navigation-only subjects here.
Neither receives rules, a correctness theorem, or a complexity claim in this
chapter.
